\documentclass[12pt]{article}
\usepackage[utf8]{inputenc}
\usepackage[T1]{fontenc}
\usepackage[margin=1in]{geometry}
\usepackage{setspace}
\usepackage{xcolor}
\usepackage[colorlinks=true, urlcolor=blue!60!black, linkcolor=blue!60!black]{hyperref}
\setlength{\parindent}{0em}
\setlength{\parskip}{0.6em}
\renewcommand{\baselinestretch}{1.0}
\begin{document}
\begin{center}
{\LARGE \textit{Texas v. Johnson}, 491 U.S. 397 (1989)\par}
\vspace{0.5em}
{\large Case Brief --- Two-Page Edition\par}
\end{center}

\section{Memory Jogger}
Burning the American flag in political protest is expressive conduct protected by the First Amendment; a State may not punish it because the message offends or to preserve the flag as a symbol of unity (\href{https://www.courtlistener.com/opinion/112304/texas-v-johnson/}{491 U.S. 397}).

\section{Detailed Case Facts}
At the 1984 Republican National Convention in Dallas, Gregory Lee Johnson burned an American flag outside City Hall at the end of a protest march while demonstrators chanted. No one was hurt, though onlookers were seriously offended. Texas convicted him of desecrating a venerated object, defined as mistreating a flag in a way the actor knows will seriously offend observers.

\section{Procedural History}
Johnson was sentenced to one year and fined \$2,000; the intermediate court affirmed, but the Texas Court of Criminal Appeals reversed. The Supreme Court affirmed.

\section{Judicial Votes}
Affirmed, 5--4. Brennan wrote for Marshall, Blackmun, Scalia, and Kennedy; Kennedy concurred; Rehnquist (with White and O'Connor) and Stevens dissented.

\section{Holding}
Flag burning in this context is expressive conduct. Because Texas's interest in preserving the flag's symbolic meaning is related to suppressing expression, \textit{O'Brien} does not apply, and the conviction fails the most exacting scrutiny.

\section{Analysis of Opinions}
Brennan found the political message ``both intentional and overwhelmingly apparent.'' No breach of the peace occurred or threatened, and offensive speech cannot be presumed to incite violence. The statute punished only offensive burning, so it turned on the message, and government may not prescribe an orthodox meaning for the flag. Kennedy acknowledged the ``personal toll'' but found the result compelled. Rehnquist likened flag burning to fighting words, ``an inarticulate grunt or roar''; Stevens saw the flag as a national asset whose protection imposes only a ``trivial burden.''

\section{Examples: Future Applications}
\textbf{Same-Side:} a federal statute banning all flag burning also fell when challenged in \textit{United States v. Eichman}. \textbf{Opposite-Side:} a protester who steals a flag or burns it in violation of a neutral fire ordinance may be prosecuted. \textbf{Fence-Sitter:} burning a flag while threatening a grieving military family and urging a crowd to attack.

\section{Critique}
Some originalists question whether burning an object is ``speech'' as originally understood. Geoffrey Stone and John Hart Ely praise the decision's enforcement of content neutrality, while critics argue it undervalues civic cohesion.

\section{Key Quotations}
\begin{quote}
``If there is a bedrock principle underlying the First Amendment, it is that the government may not prohibit the expression of an idea simply because society finds the idea itself offensive or disagreeable.''
\end{quote}
\noindent --- Justice Brennan, Opinion of the Court, 491 U.S. at 414.

\end{document}
